\section*{الفصل \ref{ch:b3:cytoskeleton-motility} --- ديناميكا الهيكل الخلوي وحركية الخلية}

\begin{solution}{exo:b3:cytoskeleton-motility:1}
الأكتين: \qty{7}{nm}، والوحدة أكتين كروي، قطبي (شائكة/مدببة)، ATP؛
وهو في القشرة والبروز والتقلص مع \omterm{def:b3:cytoskeleton-motility:motors}{الميوزين}. \omterm{def:b3:cytoskeleton-motility:filaments}{والأنيبيب الدقيق}: \qty{25}{nm}،
والوحدة ثنائي توبولين $\alpha\beta$، قطبي (موجبة/سالبة)، GTP؛ وهو في النقل
بعيد المدى والمغزل والأهداب. \omterm{def:b3:cytoskeleton-motility:filaments}{والخيط المتوسط}: \qty{10}{nm}،
وبروتينات ملفوفة لولبيًا (كيراتين، فيمنتين، لامين)، لا قطبية له ولا نوكليوتيد؛
وهو للمتانة الميكانيكية للخلية والنواة.
\end{solution}

\begin{solution}{exo:b3:cytoskeleton-motility:2}
تركيز الوحيد الذي لا تنمو عنده النهاية ولا تتقلص،
أي $k_{\text{off}}/k_{\text{on}}$. وتختلف النهايتان لأن الوحدة
تحلمه ATP فيها بعد الإدخال، فتعرض النهايتان أنواعًا
كيميائية مختلفة (أكتين يحمل ATP يصل النهاية الشائكة، وأكتين يحمل ADP
يغادر النهاية المدببة) بألفات مختلفة. ولو تساوى التركيزان
الحرجان لكان هناك توازن واحد: فلا \omterm{thm:b3:cytoskeleton-motility:treadmilling}{دوران شريطي}،
ولا تدفق، ولا سبيل إلى دفع ثمن دوران موجَّه.
\end{solution}

\begin{solution}{exo:b3:cytoskeleton-motility:3}
إلى المحيط على الأنيبيبات الدقيقة: \omterm{def:b3:cytoskeleton-motility:motors}{الكينيزين}، نحو النهاية الموجبة. وإلى المركز:
\omterm{def:b3:cytoskeleton-motility:motors}{الداينين} الهيولي، نحو النهاية السالبة. وخيط الإجهاد: \omterm{def:b3:cytoskeleton-motility:motors}{الميوزين} II،
نحو النهاية الشائكة للأكتين. \omterm{def:b3:cytoskeleton-motility:cilia}{والهدب}: \omterm{def:b3:cytoskeleton-motility:cilia}{الداينين المحوري}، نحو
النهاية السالبة للمزدوج المجاور (أي القاعدة).
\end{solution}

\begin{solution}{exo:b3:cytoskeleton-motility:4}
البروز (Rac للصفائح الزائفة، وCdc42 للأقدام الخيطية)، والالتصاق
(الإنتغرينات؛ بعد Rac وRho)، والتقلص (Rho، عبر
\omterm{def:b3:cytoskeleton-motility:motors}{الميوزين} II)، وانكماش المؤخرة (بانفكاك الالتصاقات، بدافع من
التقلص).
\end{solution}

\begin{solution}{exo:b3:cytoskeleton-motility:5}
$C_{c}^{+} = 1/5 = \qty{0.2}{\micro M}$، و$C_{c}^{-} = 0.8/1 =
\qty{0.8}{\micro M}$؛ و$C_{\text{ss}} = (1 + 0.8)/(5 + 1) = \qty{0.3}{\micro
M}$؛ و$J = 5\times 0.3 - 1 = 0.5$ وحدة في الثانية، أي \qty{1.35}{nm/s}.
\end{solution}

\begin{solution}{exo:b3:cytoskeleton-motility:6}
$800/8 = 100$ ATP في الثانية. والمتر يستغرق $1/(8\times 10^{-7}) =
1.25\times 10^{6}$ ثانية، أي نحو أسبوعين، و$1.25\times 10^{8}$ خطوة و
ATP. والانتشار كان ليستغرق \num{16000} سنة: أي فرقًا بمليار
ضعف.
\end{solution}

\begin{solution}{exo:b3:cytoskeleton-motility:7}
$F_{\max} = 8.3\times 10^{-20}/3.6\times 10^{-8} = \qty{2.3}{pN}$. فالطاقة
نفسها موزعة على خطوة أطول تعطي قوة أقل --- أي رافعة أطول.
\omterm{def:b3:cytoskeleton-motility:motors}{والكينيزين} (قصير الخطوة، \qty{6}{pN}) يناسب الأحمال الثقيلة؛ \omterm{def:b3:cytoskeleton-motility:motors}{والميوزين} V (طويل
الخطوة) يقطع مسافة أكبر لكل ATP ويناسب النقل السريع الخفيف.
\end{solution}

\begin{solution}{exo:b3:cytoskeleton-motility:8}
(أ) التاكسول يجمّد الأنيبيبات الدقيقة: فلا يستطيع المغزل أن يبحث ويلتقط و
يصحح، فيتوقف الانقسام عند نقطة التفتيش؛ والخلية الزاحفة تظل
تبرز لكنها تفقد قطبيتها الطويلة الأمد. (ب) والكولشيسين يزيل
الأنيبيبات: فلا مغزل، وتوقف انقسامي؛ ويستمر الزحف على الأكتين
لكن بلا اتجاه ثابت. (ج) والسيتوكالازين يحجب النهايات الشائكة: فلا
بروز ولا زحف؛ ويمضي الانقسام لكن يخفق انقسام الهيولى، فتنتج
خلايا ثنائية النوى. (د) وتثبيط Rac: فلا صفائح زائفة، فلا زحف؛
ولا يتأثر الانقسام كثيرًا.
\end{solution}

\begin{solution}{exo:b3:cytoskeleton-motility:9}
زمن المكث من رتبة زمن النصف، أي \qty{20}{s} (وثابت الزمن
$20/\ln 2 \approx \qty{29}{s}$)؛ والنسبة الساكنة \qty{10}{\%}.
وعلى البلمرة عند الحافة أن توفر التقدم والجريان
الراجع معًا: أي $1 + 1.5 = \qty{2.5}{\micro m/min}$.
\end{solution}

\begin{solution}{exo:b3:cytoskeleton-motility:10}
النهاية النامية تحتفظ بقلنسوة من التوبولين الحامل لجزيء GTP؛ وخلفها تنتج الحلمهة
توبولينًا حاملًا لجزيء GDP مشدودًا. وفقدُ القلنسوة --- وهو حدث عشوائي يرتفع
احتماله كلما بطؤ النمو --- يحرر الشد فتتقلص النهاية
سريعًا حتى تتكوّن قلنسوة جديدة. وفوق \omterm{thm:b3:cytoskeleton-motility:treadmilling}{التركيز الحرج}
بقليل، يكون كل أنيبيب في أي لحظة إما مقلنسًا
ناميًا وإما بلا قلنسوة منهارًا، فتتباعد الأطوال إلى جمهرة
طويلة وأخرى متلاشية بدل أن تتجمع حول متوسط. و
تكسب الخلية استكشافًا سريعًا لحجمها والقدرة على
التثبيت الانتقائي --- فالنهاية الموجبة التي تبلغ قسيمًا مركزيًا أو
موضعًا قشريًا تُلتقط وتُبقى، وتُعاد البقية تدويرًا في
دقائق: أي البحث والالتقاط.
\end{solution}

\begin{solution}{exo:b3:cytoskeleton-motility:11}
عند هذا المقياس يكون العطالة مهملًا وتكون معادلات المائع
عكوسة الزمن: فالحركة التي تعيد تتبع نفسها (مجداف يخرج
ويعود) تعيد الجسم إلى حيث بدأ، مهما كانت سرعة كل
ضربة. أما اللولب الدوّار فلا يعيد تتبع نفسه قط --- فحركته
كيرالية متصلة --- فينتج دفعًا صافيًا. وعند \qty{100}{Hz}
مع ألف بروتون في الدورة، يكون $10^{5}$ بروتون في الثانية لكل
محرّك.
\end{solution}

\begin{solution}{exo:b3:cytoskeleton-motility:12}
لا تستطيع أهداب الطرق الهوائية أن تضرب، فلا يُنظَّف المخاط ولا البكتيريا،
وتتكرر الأخماج. وأسواط النطاف، المبنية على \omterm{def:b3:cytoskeleton-motility:cilia}{المحور الهدبي} نفسه،
لا تتحرك: فالعقم. وفي الجنين تقود أهداب العقدة المتحركة
جريانًا نحو اليسار يضبط المحور الأيسر--الأيمن؛ ومن غير جريان يُختار الجانب
عشوائيًا، فيكون نصف المرضى معكوسي الأحشاء.
\end{solution}

\begin{solution}{pb:b3:cytoskeleton-motility:1}
\textbf{1.} $\qty{300}{\micro m^3} = 3\times 10^{-13}$ لتر؛ $\times
2\times 10^{-4}$ مول/لتر $= 6\times 10^{-17}$ مول $= 3.6\times 10^{7}$
جزيء؛ منها $1.8\times 10^{7}$ في الخيوط.
\textbf{2.} $1.8\times 10^{7}\times \qty{2.7}{nm} = \qty{4.9}{cm}$ من
الخيوط؛ وعند \qty{0.25}{\micro m} لكل خيط، نحو $2\times 10^{5}$
خيط.
\textbf{3.} بروتين الحجب يحجب معظم النهايات الشائكة، ويربط البروفيلين و
الثيموزين-$\beta$4 الوحيدات حتى يكون الأكتين الحر حقًا الحامل لجزيء ATP
قرب \omterm{thm:b3:cytoskeleton-motility:treadmilling}{التركيز الحرج}؛ فينحصر النمو في النهايات التي
تفكّ الخلية حجبها.
\textbf{4.} $C_{\text{ss}} = 2.2/12.9 = \qty{0.17}{\micro M}$؛ و$J = 11.6
\times 0.17 - 1.4 = 0.6$ وحدة في الثانية؛ أي \qty{1.6}{nm/s} $=
\qty{0.1}{\micro m/min}$.
\textbf{5.} $1/0.1 = 10$ أضعاف.
\textbf{6.} \qty{1}{\micro m/min} $= \qty{16.7}{nm/s} = 6.2$ وحدة في
الثانية لكل خيط؛ ومضروبًا في $\times 2\times 10^{5} = 1.2\times 10^{6}$ ATP في
الثانية، أي نحو $0.1\,\%$ من دوران الخلية.
\textbf{7.} $1/(8\times 10^{-7}) = 1.25\times 10^{6}$ ثانية $\approx 14.5$
يومًا؛ و$1.25\times 10^{8}$ خطوة وATP.
\textbf{8.} $L^{2}/2D$: أي $(1)^{2}/(2\times 10^{-12}) = 5\times 10^{11}$ ثانية
$\approx \num{16000}$ سنة للمتر؛ و$(10^{-5})^{2}/(2\times
10^{-12}) = 50$ ثانية لمسافة \qty{10}{\micro m}. فالانتشار يفي داخل
جسم الخلية؛ وما زاد على ذلك يحتاج إلى محركات.
\textbf{9.} $5\times 10^{4}/6.02\times 10^{23} = 8.3\times 10^{-20}$ جول
$= 20\,k_{B}T$.
\textbf{10.} $8.3\times 10^{-20}/8\times 10^{-9} = \qty{10}{pN}$؛
والكفاءة $6/10 \approx 60\,\%$.
\textbf{11.} $F = 6\pi\times 0.01\times 10^{-7}\times 8\times 10^{-7} =
1.5\times 10^{-14}$ نيوتن $= \qty{0.015}{pN}$، أي أقل بأربعمئة مرة من
\omterm{prop:b3:cytoskeleton-motility:physics}{قوة التوقف}: فمحرّك واحد يكفي ويزيد في مواجهة المقاومة اللزجة؛ والأحمال
الحقيقية هي العوائق والأربطة.
\textbf{12.} $10^{4}\times 100 = 10^{6}$ ATP في الثانية، أي $0.1\,\%$ من
ميزانية العصبون: فالنقل رخيص. لكن المحركات تتوقف لحظة يهبط ATP،
فإخفاق المتقدرات يجوّع المشبك من كل ما يُصنع في جسم
الخلية --- وهو إخفاق النقل المحوري المرصود في
أمراض التنكس العصبي.
\textbf{13.} \qty{10}{\micro m/min} $= \qty{167}{nm/s}$؛ و$167/2.7 = 62$
وحدة في الثانية لكل خيط.
\textbf{14.} $200\times 20 = 4000$ خيط؛ و$4000\times 62 = 2.5\times
10^{5}$ وحدة في الثانية.
\textbf{15.} $2.5\times 10^{5}$ ATP في الثانية، أي $0.025\,\%$ من
الميزانية.
\textbf{16.} $\qty{3}{\micro m}/(\qty{10}{\micro m/min}) = \qty{0.3}{min}
= \qty{18}{s}$؛ و$2.5\times 10^{5}\times 18 = 4.5\times 10^{6}$ وحدة
في الشبكة.
\textbf{17.} $1\,\text{pN}\times 2.7\,\text{nm} = \qty{2.7}{pN.nm} =
0.66\,k_{B}T$: فتقلب بهذه الطاقة يقع باحتمال
نحو $e^{-0.66} \approx 0.5$ --- أي تنفتح ثغرات بحجم وحدة واحدة باستمرار،
والسقّاطة معقولة تمامًا.
\textbf{18.} البلمرة نفسها تدفع، والبكتيريا هي
``الغشاء'' في مقدمة شبكتها هي؛ ومع آلة Arp2/3 عند
الخلية تبلغ سرعة \omterm{def:b3:cytoskeleton-motility:migration}{الصفيحة الزائفة} وأكثر --- أي ما يصل إلى
\qty{0.5}{\micro m/s}، أو \qty{30}{\micro m/min}.
\textbf{19.} عند $C = K_{d}$ يكون $\theta = 1/2$ و$\mathrm{d}\theta/
\mathrm{d}C = 1/(4C)$: ففرق \qty{2}{\%} في $C$ يعطي $\Delta\theta
= 0.005$؛ ومع \num{25000} مستقبل لكل جانب، يكون $125$ مشغولًا زيادةً في
المقدمة.
\textbf{20.} لكل جانب نحو \num{12500} مشغول، يتقلب بمقدار
$\sqrt{\num{12500}} \approx 110$، فيتقلب فرق الجانبين
بنحو $160$: أي يضيع العدد $125$ في الضجيج في أي
لحظة. والمستقبلات تعيد الارتباط نحو مرة في الثانية؛ والتوسيط على مدى $N$
ثانية يقلّص الضجيج بمقدار $\sqrt{N}$ --- فعشر ثوانٍ تنزل به إلى
$50$ ويُقرأ التدرج.
\textbf{21.} Rac (مع Cdc42) في المقدمة، وRho في المؤخرة. ومع Rac
فعّال في كل مكان تبرز الخلية من كل جانب، وتنبسط، ولا تذهب إلى شيء.
\textbf{22.} $30/3 = \qty{10}{\micro m/min}$: كما أُعطي.
\textbf{23.} يُستبدل بالبروز نفطٌ يقوده الضغط؛ أما الالتصاق
والتقلص فما زالا يتوقفان على الأكتين \omterm{def:b3:cytoskeleton-motility:motors}{والميوزين}، فيبقى الأكتين
جوهريًا لبقية الدورة.
\textbf{24.} تُعاد بناء المقدمة كل \qty{18}{s}: فإزاحة في نشاط Rac
تعيد توجيه البلمرة خلال عمر شبكة واحد، فتقود المقدمة
الجديدة، وتتبعها الخطوات البطيئة في المؤخرة.
\textbf{25.} \omterm{thm:b3:cytoskeleton-motility:treadmilling}{الدوران الشريطي} خارج الجسم \qty{1.6}{nm/s}؛ ونحو $14.5$ يومًا
لعبور المحور؛ ونحو $2.5\times 10^{5}$ ATP في الثانية لأجل
البروز.
\end{solution}
